% CS 418: Intro to Data Science — Worksheet 24 (24-week12-day2)
% Week 12, Wednesday 2026-11-11
% Compile with: latexmk -pdf worksheet.tex

\documentclass[11pt]{article}

% ── Packages ──────────────────────────────────────────────────────────────────
% The wide right margin is the open-ended notes column, half of the 7.2in between the outer margins less a 0.3in gutter; the header and title block span both.
\usepackage[margin=0.65in, right=4.4in, headheight=14pt]{geometry}
\usepackage{amsmath, amssymb}
\usepackage{ragged2e}
% Never break a word across lines, including at hyphens ("Sans-Serif").
\hyphenpenalty=10000
\exhyphenpenalty=10000
\usepackage{enumitem}
\usepackage{booktabs}
\usepackage{graphicx}
\usepackage{hyperref}
\usepackage{xcolor}
\usepackage{fancyhdr}
\usepackage{mdframed}
\usepackage{array}
\usepackage{tabularx}
\usepackage{tikz}

% ── Draft watermark ───────────────────────────────────────────────────────────
% This worksheet is still a stub. The watermark marks the PDF as unfinished so
% nobody mistakes a placeholder for the real thing. Delete this block once the
% worksheet has real content.
\usepackage{draftwatermark}
\SetWatermarkText{IN-PROGRESS}
\SetWatermarkScale{2.2}
\SetWatermarkLightness{0.85}

% ── Page style ────────────────────────────────────────────────────────────────
\pagestyle{fancy}
\fancyhf{}
\fancyhfoffset[R]{3.75in}
\lhead{\textbf{CS 418: Intro to Data Science}}
\rhead{UIC, Fall 2026}
\cfoot{\thepage}
\renewcommand{\headrulewidth}{1pt}
\setlength{\headheight}{12pt}
\setlength{\headsep}{0.4cm}

% ── Hyperlinks ────────────────────────────────────────────────────────────────
\hypersetup{colorlinks=false, hidelinks}

% ── Convenience environments ──────────────────────────────────────────────────
% Usage: \fixedblock[height]{content} — a full-width block exactly that tall.
\newcommand{\fixedblock}[2][5.7cm]{\noindent\begin{minipage}[t][#1][t]{\linewidth}\raggedright #2\end{minipage}\par}

% Usage: \vocabrow{term} — one row of checkpoint 1's one-column vocab table: the term in bold with blank space below it for the definition.
\newcommand{\vocabrow}[1]{\small \textbf{#1}\rule[-1.5em]{0pt}{0.1pt}}

% ══════════════════════════════════════════════════════════════════════════════
\begin{document}
\RaggedRight

% ── Title block: topic left, info box right ───────────────────────────────────
% Spans the notes column too, through the column headings; \makebox lets it run 3.75in past \textwidth without an overfull warning.
\noindent\makebox[\textwidth][l]{\begin{minipage}[t]{\dimexpr\textwidth+3.75in\relax}
\vspace{0pt}
\begin{minipage}[t]{0.55\linewidth}
  \vspace{0pt}
  {\LARGE \textbf{Worksheet \#\,24}} \\[0.18em]
  {\large \textit{Week 12, Day 2: Network immunization}}
\end{minipage}
\hfill
\begin{minipage}[t]{0.40\linewidth}
  \vspace{0pt}
  \small
  Name: \underline{\hspace{4.5cm}} \\[0.35em]
  NetID: \underline{\hspace{4.3cm}} \\[0.35em]
  Date: \underline{\hspace{4.4cm}} \\
\end{minipage}

\vspace{0.6em}
\noindent\rule{\linewidth}{0.6pt}
\vspace{0.1em}

\noindent
\begin{minipage}[t]{3.45in}
  \noindent\underline{\textbf{Guided checkpoints}}\\[0.2em]
  \noindent{\small For full credit, fill in the checkpoints. If you skip a checkpoint, include your own notes for full credit.}
\end{minipage}%
\hspace{0.3in}%
\begin{minipage}[t]{3.45in}
  \noindent\underline{\textbf{Open-ended notes}}\\[0.2em]
\end{minipage}
\end{minipage}}

\vspace{0.8em}

% ── 1. Vocab ─────────────────────────────────────────────────────────────────
% Key: definitions from the Day 2 vocab slides (MMDS Ch. 10; quotes mostly from the summary, Sec. 10.9), in deck order.
\fixedblock[0.36\textheight]{%
\noindent\textbf{1. Vocab}\\[0.18em]
\noindent Write a short definition and/or example for each.

\vspace{0.5em}
\noindent
{\renewcommand{\arraystretch}{1.15}%
\begin{tabular}{|l|}
  \hline
  \vocabrow{Community; clique}\\ \hline
  \vocabrow{Adjacency matrix}\\ \hline
  \vocabrow{Degree matrix}\\ \hline
  \vocabrow{Overlapping communities}\\ \hline
  \vocabrow{Neighborhood of radius $r$; diameter}\\ \hline
  \vocabrow{Transitive closure}\\ \hline
\end{tabular}}
}

\end{document}
